\section{Luận đề của cuốn sách}
\label{sec:ch01-luan-de-cua-sach}

Luận đề của cuốn sách không phải là mọi giải thích đều vô ích, cũng không
phải là mô hình phức tạp luôn phải bị thay bằng mô hình đơn giản. Luận đề hẹp
hơn thế: một thước đo được gọi là \tn{độ trung thực}{faithfulness} phải được
kiểm chứng như một dụng cụ đo. Nếu chưa biết điểm số của nó thay đổi ra sao
khi cơ chế thật thay đổi, hoặc nó có phân biệt được lời giải thích đúng cơ
chế với lời giải thích chỉ nghe hợp lý hay không, thì ta chưa có lý do để xem
điểm số ấy là bằng chứng về độ trung thực.

Các chương kế tiếp đi theo thứ tự cần thiết cho lập luận đó.
Chương~\ref{ch:mo-hinh-can-giai-thich} thu gọn nền tảng về các mô hình mà
những lời giải thích sau này sẽ nhắm tới. Các chương về LIME, SHAP và
attribution theo gradient trình bày cơ chế, giả định và đầu ra của từng họ
phương pháp. Sau khi cơ chế, giả định và đầu ra của từng họ đã rõ,
chương~\ref{ch:do-do-trung-thuc} phân loại các họ chỉ số độ trung thực, còn
chương~\ref{ch:chi-so-khong-do-do-trung-thuc} kiểm tra trực diện tuyên bố
rằng các chỉ số ấy đo được độ trung thực.

Trình tự này cũng nhằm tránh một cách đọc sai phổ biến: khi một bản đồ nhiệt
hoặc danh sách đặc trưng trông thuyết phục, ta dễ nhảy thẳng từ hình thức của
lời giải thích sang kết luận về mô hình. Cuốn sách giữ hai bước tách biệt:
trước hết mô tả phương pháp đang tạo ra cái gì, sau đó mới hỏi bằng chứng nào
cho thấy nó bám vào hành vi mô hình. Câu hỏi sau khó hơn, và chính nó quyết
định ta được phép đặt bao nhiêu niềm tin vào lời giải thích.
